\section{Mécanique relativiste}

    \subsection{Énergie et impulsion relativistes}

        \subsubsection{Quadri-vecteur énergie-impulsion}

        \begin{omed}{Définition (quadri-vecteur impulsion)}
            On appelle \textbf{quadri-vecteur impulsion} ou \textbf{impulsion} le quadri-vecteur 
            \[ P = m U = \left(\frac{E}{c}, \mathbf{p}\right) \] 
            où  
            \begin{align*}
                E &= \gamma m c^2 = \frac{mc^2}{\sqrt{1 - \frac{v^2}{c^2}}} \\
                \mathbf{p} &= \gamma m \mathbf{v} = \frac{m \mathbf{v}}{\sqrt{1 - \frac{v^2}{c^2}}}
            \end{align*}
            sont l’\textbf{énergie totale relativiste} et l’\textbf{impulsion relativiste}.
        \end{omed}

        \subsubsection{Particules de masse nulle}

    \subsection{relation fondamentale de la dynamique}

    \begin{omed}{Relation fondamentale de la dynamique}
        En relativité, la relation fondamentale de la dynamique s’écrit
        \[ \frac{\text{d}P}{\text{d}\uptau} = \mathcal{F} \]   
        où $\mathcal{F}$ est le \textbf{quadri-vecteur force de Minkowski} ou \textbf{quadri-force}. 
    \end{omed}

    Dans le cas particulier d’une particule massive, cette relation peut se réécrire 
    \[ m \mathcal{A} = \mathcal{F} \]   
    La propriété $U \cdot A = 0$ implique que
    \[ U \cdot \mathcal{F} = 0 \]   
    Dans le référéntiel instantané $\mathcal{R}_*$ où $U_* = \left(c, \vec{0}\right)$, la composante temporelle de $\mathcal{F}$ est donc nulle.

        